\section{Preliminaries}
\label{preliminaries}

We fill in some details on network structure and penetration
testing. We give a brief background on POMDPs.



\subsection{Network Structure}


Networks can be viewed as directed graphs whose vertices are given by
the set $M$ of \emph{machines}, and whose arcs are connections between
pairs of $m \in M$. However, in practice, these network graphs have a
particular structure. They tend to consist of \emph{subnetworks}, \ie,
clusters $N$ of machines where every $m \in N$ is directly connected
to every $m' \in N$. %
By contrast, not every subnetwork $N$ is connected to every other
subnetwork $N'$, and typically, if such a connection does exist, then
it is filtered by a \emph{firewall}.

From the perspective of an attacker, the firewalls filter the
connections and thus limit the attacks that can be executed when
trying to hack into a subnetwork $N'$ from another subnetwork $N$.  On
the other hand, once the hacker managed to get into a subnetwork $N$,
access to all machines within $N$ is easy. Thus a natural
representation of the network, from an attack planning point of view,
is that of a graph whose vertices are subnetworks, and whose arcs are
annotated with firewalls $F$. We herein refer to this graph as the
\emph{logical network} $LN$, and we denote its arcs with $N
\xrightarrow{F} N'$.

We formalize firewalls as sets of rules describing which kinds of
communication (\eg, ports) are disallowed. Thus smaller sets
correspond to ``weaker'' firewalls, and the \emph{empty firewall}
blocks no communication at all.


We remark that, in our POMDP model, we do not provide for privilege
escalation, or obtaining passwords. This can instead be modeled at the
level of $LN$. Different privilege levels on the same machine $m$ can
be encoded via different copies of $m$. If controlling $m$ allows
the retrieval of passwords, then $m$ can be connected via empty firewalls
to the machines $m'$ who can be accessed by using these passwords, more
precisely to high-privilege copies of these $m'$.





\subsection{Penetration Testing}


Uncertainty in pentesting arises because it is impossible to keep
track of all the \emph{configuration} details of individual machines,
i.e., exactly which versions of which programs are installed
etc. However, it is safe to assume that the pentesting tool knows the
structure of the network, \ie, the graph $LN$ and the filtering done
by each firewall: changes to this are infrequent and can easily be
registered.


The objective of pentesting is to gain control over certain machines
(with critical content) in the network. At any point in time, each
machine has a unique {\em status}. A \emph{controlled} machine $m$ has
already been hacked into. A \emph{reached} machine $m$ is connected to
a controlled machine, \ie, either $m$ is in a subnetwork $N$ one of
whose machines is controlled, or $m$ is in a subnetwork $N'$ with a
$LN$ arc $N \xrightarrow{F} N'$ where one of the machines in $N$ is
controlled. All other machines are \emph{not reached}. The algorithm
starts with one controlled machine, denoted here by
\start.\footnote{For simplicity, we will notate \start\ as a separate
  vertex in $LN$.  If \start\ is part of a subnetwork $N$, this means
  to turn $N \setminus \{\start\}$ into a separate vertex in $LN$,
  connected to \start\ via the empty firewall.}
We will use the following (small but real-life) situation as a running
example:

\vspace{-0.05cm}
\begin{example}\label{exp:running-preliminaries}
The attacker has already hacked into a machine $m'$, and now wishes to
attack a machine $m$ within the same subnetwork. The attacker knows
two exploits: \emph{SA}, the ``Symantec Rtvscan buffer overflow
exploit''; and \emph{CAU}, the ``CA Unicenter message queuing
exploit''. SA targets a particular version of ``Symantec Antivirus'',
that usually listens on port 2967. CAU targets a particular version of
``CA Unicenter'', that usually listens on port 6668. Both work only if
a protection mechanism called \emph{DEP} (``Data Execution
Prevention'') is disabled.
\end{example}
\vspace{-0.05cm}

If SA fails, then it is likely that CAU will fail as well (because DEP
is enabled). The attacker is then better off trying something
else. Achieving such behavior requires the attack plan to observe the
outcomes of actions, and to react accordingly. Classical planning
(which assumes perfect world knowledge at planning time) cannot
accomplish this.  

Furthermore, port scans---observation actions testing whether or not a
particular port is open---should be used only if one actually intends
to execute a relevant exploit. Here, if we start with SA, we should
scan only port 2967. We accomplish such behavior through the use of
POMDPs. By contrast, to reduce uncertainty, classical planning
requires a pre-process executing \emph{all} possible scans. In this
example, there are only two---ports 2967 and 6668---however in general
there are many, causing significant network traffic and waiting time.





\subsection{POMDPs}


POMDPs are usually defined (e.g., \cite{Monahan82,KaeLitCas-aij98}) by
a tuple $\langle \cS, \cA, \cO, T, O, r, b_0 \rangle$. If the system
is in state $s \in \cS$ (the {\em state space}), and the agent
performs an action $a\in \cA$ (the {\em action space}), then that
results in (1) a transition to a state $s'$ according to the {\em
  transition function} $T(s,a,s')=Pr(s'| s,a)$, (2) an observation
$o\in\cO$ (the {\em observation space}) according to the {\em
  observation function} $O(s',a,o)=Pr(o| s',a)$ and (3) a scalar {\em
  reward} $r(s,a,s')$. $b_0$, the {\em initial belief}, is a
probability distribution over $\cS$. %

The agent must find a decision {\em policy} $\pi$ choosing, at each
step, the best action based on its past observations and actions so as
to maximize its future gain, which we measure here through the total
accumulated reward. The expected value of an optimal policy is denoted
with $V^*$.


The agent typically reasons about the hidden state of the system using
a {\em belief state} $b$, a probability distribution over $\cS$. For
our experiments we use SARSOP \cite{KurHsuLee08}, a state of the art
point-based algorithm, i.e., an algorithm approximating the value
function as the upper envelope of a set of hyperplanes, corresponding
to a selection of particular belief states (referred to as
``points'').



